\section{System Model and Adversary Model}
\label{sec:sysmodel}

\subsection{System Model}

PhysCausal involves two principals: a \textit{client device} (the smartphone under
test) and a \textit{detection server} operated by a platform or application provider
(e.g., a mobile bank, ad exchange, or app store).%

\noindent\textbf{Client device.}
The client runs a lightweight PhysCausal agent embedded in the host application.%
The agent operates in one of two server-selected modes.%

In \textit{passive mode}, the agent performs two concurrent actions during the
monitoring window: (i)~records \texttt{TYPE\_LINEAR\_ACCELERATION} at
\texttt{SENSOR\_DELAY\_FASTEST} and (ii)~logs touch-event timestamps reported by the
Android input system.%
No interaction is requested from the user and no visible UI change occurs.%

In \textit{active mode}, the server additionally selects a challenge time
$t^* \in [\frac{T}{4}, \frac{3T}{4}]$ \textit{after} the monitoring window begins and
notifies the app.%
The app presents a natural UI element (e.g., a tappable button or confirmation widget)
at $t^*$.%
The agent records the Android-reported touch timestamp and transmits both streams to
the server.%
Crucially, $t^*$ is communicated only after recording starts, so any pre-recorded
accelerometer stream injected by an adversary cannot contain a correctly timed
impulse transient.%

No inference runs on-device; the full PCCN pipeline executes server-side, keeping
the detection model confidential and preventing white-box adaptive attacks.%

\noindent\textbf{Detection server.}
The server receives the encrypted payload, extracts the linear acceleration signal and
touch-event timestamps, and runs PCCN in passive or active mode.%
In passive mode, PCCN classifies using $\delta(t)$ and $\tau(t)$ alone.%
In active mode, PCCN additionally verifies the causal coupling at $t^*$.%
The server returns a ternary verdict (\textit{human}, \textit{autonomous agent},
or \textit{overlay agent}) to the calling service for downstream enforcement.%

\noindent\textbf{Deployment properties.}
The agent requires no dangerous Android permissions.%
Passive monitoring adds zero user friction.%
Active mode presents a single natural UI element indistinguishable from normal app flow.%
Detection is triggered server-side, so a farm operator cannot anticipate or selectively
suppress individual checks.%

\textcolor{green}{[Figure: System architecture diagram. Left column (Client Device):
smartphone icon at top; three stacked blocks connected by arrows:
(1) ``Linear Acceleration Recording'' (gray, waveform icon, zero-permission label),
(2) ``Touch Event Logging'' (gray, tap-finger icon),
(3) ``Active Mode: Challenge at $t^*$'' (gray, dashed border, clock icon, label
``$t^*$ chosen after recording starts'').
Arrow to right: ``Encrypted Payload.''
Right column (Detection Server): three stacked blocks:
(1) ``Baseline Estimation $\rightarrow$ $\delta(t)$'' (blue),
(2) ``PCCN Classifier'' split into ``Passive: $\delta(t) + \tau(t)$'' and
``Active: $\delta(t) + \tau(t)$ at $t^*$'' (blue, dashed divider),
(3) output diamond ``Human / Autonomous Agent / Overlay Agent'' (amber/navy/gray tricolor).
Arrow from output: ``Verdict to App Backend.''
White background, unified color palette.]}

\subsection{Adversary Model}
\label{sec:adversary}

The adversary is an LLM agent operator whose goal is to cause agent-controlled sessions
to be labeled as human-operated.%
We consider four adversary levels of increasing capability.%

\noindent\textbf{A0 (Default --- no evasion).}
The agent dispatches touch events via ADB or \texttt{AccessibilityService} with no
attempt to modify sensor output.%
The accelerometer shows no impulse at any touch-event timestamp.%
This is trivially detected by PhysCausal.%

\noindent\textbf{A1 (Vibration motor).}
The adversary attaches a vibration motor to the phone chassis and triggers it
synchronously with each injected touch event, producing a mechanical signal in the
accelerometer.%
Motor-induced vibrations have a different spectral signature from real finger-impact
impulses: motors produce narrowband sinusoidal oscillation at the motor's driving
frequency, whereas finger impact produces a broadband transient with characteristic
exponential decay.%
In active mode, the vibration motor must additionally fire at the server-selected $t^*$,
which is communicated only after recording begins; the motor controller can comply, but
the spectral and temporal waveform mismatch remains detectable.%

\noindent\textbf{A2 (Timing mimicry).}
The adversary adjusts inter-touch intervals to match the lognormal distribution of
human timing, defeating the timing-distribution stream of PCCN.%
Active mode remains fully effective: mimicking timing does not produce physical impulses.%
Passive mode retains residual detection from the causality stream alone.%

\noindent\textbf{A3 (Human operator).}
The adversary assigns a human to physically hold and tap the phone for each session,
producing authentic physical signals.%
This eliminates the economic advantage of LLM automation (1:1 human-to-session ratio)
and is treated as outside the practical threat model.%
